\documentclass[12pt]{article}
\usepackage{amsmath}
\begin{document}
\textit{Proposed proposed to SSMJ }\\ \\
\textit{Problem proposed by Prakash Pant, The University of Vermont, Bardiya, Nepal.}\\ \\
\textbf{Matrix Equation
} \\ \\
\textbf{Statement of the Problem:}
Find the \( 2 \times 2\) matrices M with real entries that satisfy the equation:\\
\[ M^2 - 5M  = 
\begin{pmatrix}
0 & 9 \\
4 & 0
\end{pmatrix}
\] 
\textbf{Solution of the Problem:} \\

Notice that \[ M^2 - 5M + 6I = \begin{pmatrix} 6 & 9 \\ 4 & 6 \end{pmatrix} \]
\[ (M - 2I) ( M - 3I) = \begin{pmatrix} 6 & 9 \\ 4 & 6   \end{pmatrix} \]
Taking determinant on both sides 
\[ (|M-2I|)\,\, (|M-3I|) = 0 
\]
\[ |M-2I| = 0 \text{ or } |M-3I| = 0 \]
This means, if \( \lambda_1\) be one of the eigenvalues of the matrix M, then 
\begin{equation} \lambda_1 = 2 \text{ or } 3 \end{equation}
Similarly, also notice that 
\[ M^2 - 5M -6I = \begin{pmatrix}  -6 & 9 \\ 4 & -6  \end{pmatrix} \]
\[ (M+I) (M-6I) = \begin{pmatrix} -6 & 9 \\ 4 & -6   \end{pmatrix} \]
Taking determinant on both sides 
\[ (|M+I|) (|M-6I|) = 0 \]
\[ |M+I|=0 \text{ or } |M-6I| = 0 \]
This means, if \( \lambda_2\) be the other eigenvalue of the matrix M, then 
\begin{equation} \lambda_2 = -1 \text{ or } 6 \end{equation}
We won't explore any further for eigenvalues because a \( 2 \times 2\) Matrix can only have two eigenvalues. \\
We have four possible pairs of eigenvalues for matrix M using equation (1) and (2). 
\[ (\lambda_1,\lambda_2) = (2,-1) \text{ or } (2,6) \text{ or } (3,-1) \text{ or } (3,6)\]
We know, for a \(2 \times 2\) Matrix, the characteristic equation is : 
\[ \lambda^2 - \text{Tr(M)} \lambda + \text{Det(M)} = 0 \]
where Tr(M) stands for Trace of a Matrix and Det(M) stands for Determinant of a Matrix.\\
By Cayley-Hamilton theorem, any matrix M satisfies its characteristic equation. Thus, 
\[M^2 - \text{Tr(M)} M + \text{Det(M)} = 0 \]
Using the fact that Tr(M) = \( \lambda_1 + \lambda_2\) and Det(M) = \( \lambda_1 \lambda_2\), \\
For eigenvalue pair (2,-1), we get \begin{equation} M^2 - M - 2I = 0 \implies M^2 = M + 2I \end{equation}
which is a necessary but not sufficient equation.  Using the Problem statement and applying equation (3) on it
\[ M^2 - 5M = \begin{pmatrix} 0 & 9 \\ 4 & 0 \end{pmatrix} \]
\[ M + 2I - 5M  = \begin{pmatrix} 0 & 9 \\ 4 & 0 \end{pmatrix}  \implies  -4M + 2I = \begin{pmatrix} 0 & 9 \\ 4 & 0 \end{pmatrix} \]

\[\implies  -4M = \begin{pmatrix} -2 & 9 \\ 4 & -2 \end{pmatrix} \implies M = \begin{pmatrix} \frac{1}{2}
& -\frac{9}{4} \\ -1 & \frac{1}{2} \end{pmatrix} \]
Repeating the same procedure, for eigenvalue pair (2,6), we get \( M = \begin{pmatrix}4 & 3 \\ \frac{4}{3} & 4 \end{pmatrix} \). For eigenvalue pair (3,-1), we get \(M = \begin{pmatrix}1 & -3 \\ -\frac{4}{3} & 1 \end{pmatrix} \). For eigenvalue pair (3,6), we get \( M = \begin{pmatrix}\frac{9}{2} & \frac{9}{4} \\ 1 & \frac{9}{2} \end{pmatrix} \). These four solutions are all possible solutions for real matrices M. 









\end{document}